\lyseis{1}{Επανάληψη ενότητας 1}
\setcounter{chapter}{0}
\solr{1}{1} $37.002.005$. Στα εκατομμύρια: το ψηφίο των εκατοντάδων χιλιάδων είναι 0, άρα $37.000.000$.
\solr{1}{2} $4.999,\ 5.012,\ 5.021,\ 5.102,\ 10.000$.
\solr{1}{3} Δεκάδες: $86.550$ (μονάδες 9). Εκατοντάδες: $86.500$ (δεκάδες 4). Χιλιάδες: $87.000$ (εκατοντάδες 5).
\solr{1}{4} α) $25\cdot 4\cdot 17=1.700$ \quad β) $99\cdot 8=800-8=792$ \quad γ) $36\cdot(15+85)=3.600$ \quad δ) $(47+53)+(68+32)=200$.
\solr{1}{5} α) $3\cdot 36-10=98$ \quad β) $32-25=7$ \quad γ) $20-8=12$ \quad δ) $100-4\cdot 25=0$.
\solr{1}{6} $1.000=48\cdot 20+40$. Ο αριθμός είναι $15\cdot 7+6=111$.
\solr{1}{7} $100=7\cdot 14+2$: Πέμπτη και 2 μέρες, δηλαδή Σάββατο.
\solr{1}{8} Με 2 (τελειώνει σε 0), 5, 10, 3 και 9 (άθροισμα ψηφίων 9), 4 (40 διαιρείται με το 4). Όχι με το 25 (40 δεν είναι 00, 25, 50, 75).
\solr{1}{9} $5+7=12$, και χρειάζεται 18, άρα $\square=6$: $567=9\cdot 63$.
\solr{1}{10} Πρώτοι: 59 και 67. Σύνθετοι: $57=3\cdot 19$ και $63=3^2\cdot 7$.
\solr{1}{11} $360=2^3\cdot 3^2\cdot 5$. Έχει 3 διαφορετικούς πρώτους παράγοντες.
\solr{1}{12} $36=2^2\cdot 3^2$ και $60=2^2\cdot 3\cdot 5$: $\MKD=2^2\cdot 3=12$, $\EKP=2^2\cdot 3^2\cdot 5=180$.
\solr{1}{13} Ο αριθμός των στοιβών είναι ο $\MKD(36,\,60)=12$. Σε κάθε στοίβα: $36:12=3$ κάρτες με ζώα και $60:12=5$ με αυτοκίνητα.
\solr{1}{14} $\EKP(8,\,12)=24$: σε 24 μέρες.
\solr{1}{15} α) Λάθος (το 6 είναι άρτιος και δεν διαιρείται με το 4). β) Σωστό. γ) Λάθος: $91=7\cdot 13$. δ) Λάθος: $2^4=16$. ε) Λάθος: $17-(9-3)=11$, ενώ $17-9-3=5$. στ) Σωστό.
\solr{1}{16} Δίκιο έχει η Ζωή: ο πολλαπλασιασμός γίνεται πρώτος, $4\cdot 5=20$, και $6+20=26$. Για να βγει 50: $(6+4)\cdot 5=50$.
\solr{1}{17} $100=2^2\cdot 5^2$. Οι διαιρέτες είναι 1, 2, 4, 5, 10, 20, 25, 50, 100: 9 διαιρέτες. Πρώτοι είναι οι 2 και 5, δηλαδή 2.
